\chapter{Implementacja i architektura pakietu}

% Odpowiednik sekcji Software Description. Najważniejszy rozdział inżynierski:
% opisujesz własny wkład programistyczny. Objętość: 10-14 stron.

Pakiet zaimplementowano w języku R zgodnie ze standardami repozytorium CRAN. Kod źródłowy jest otwarty i dostępny w repozytorium podanym na stronie tytułowej.

\section{Struktura projektu}

Opisz katalogi i zależności. Wyjaśnij, co gdzie leży i dlaczego. Wymień zależności z podziałem na te niezbędne do działania i te potrzebne wyłącznie do zbudowania dokumentacji i testów.

\begin{tabelaepi}{Moduły pakietu i ich odpowiedzialność}{tab:moduly}
\begin{tabular}{@{}lll@{}}
\toprule
Plik & Odpowiedzialność & Funkcje eksportowane \\
\midrule
\plik{przygotowanie\_danych.R} & walidacja i normalizacja & \fun{przygotuj\_dane} \\
\plik{algorytm.R}              & silnik obliczeniowy      & \fun{oblicz\_wynik} \\
\bottomrule
\end{tabular}
\zrodlo{oprac. własne}
\end{tabelaepi}

\section{Przetwarzanie i walidacja danych}

Opisz drogę od danych surowych do struktury, na której pracuje algorytm: sprawdzanie warunków wstępnych, obsługę braków i kodów błędnych, skalowanie, agregację. Wyjaśnij, które naruszenia zatrzymują wykonanie i z jakim komunikatem – to jest kontrakt Twojego narzędzia.

\section{Silnik obliczeniowy i klasy wyniku}

Wyjaśnij, jak wzory z rozdziału pierwszego przełożyły się na kod. Pokaż fragment kodu, nie całość – całość jest w repozytorium.

% Wnętrze listingu składa się wyłącznie ze znaków ASCII. Komentarze w kodzie
% pisz bez znaków diakrytycznych. Sprawdzenie: Rscript tools/sprawdz-listingi.R
\begin{lstlisting}[caption={Wyznaczanie wag metoda entropii},label={lst:wagi}]
oblicz_wagi_entropia <- function(macierz) {
  stopifnot(is.matrix(macierz), all(macierz >= 0))

  p <- sweep(macierz, 2, colSums(macierz), "/")
  e <- -colSums(p * log(p + .Machine$double.eps)) / log(nrow(macierz))
  d <- 1 - e

  d / sum(d)
}
\end{lstlisting}

Opisz zastosowany system klas i metody generyczne. Wyjaśnij, dlaczego wynik jest obiektem określonej klasy, a nie zwykłą ramką danych.

\section{Procedura generowania danych testowych}

Opisz procedurę generowania danych o znanej strukturze: rozkłady, parametry, zależności między zmiennymi, kontrolowane braki i wartości odstające, ziarno losowości. To jest element metodologii, nie dodatek – dzięki niemu wiadomo, jaki wynik narzędzie powinno zwrócić.

\section{Weryfikacja poprawności}

Opisz trzy warstwy sprawdzania: przypadki analityczne o wyniku policzonym ręcznie, testy własności (niezmienniczość na skalowanie, monotoniczność, permutacja etykiet) oraz testy odporności na dane niepoprawne. Podaj wynik sprawdzenia zgodności z wymaganiami repozytorium na trzech systemach operacyjnych.

\section{Wykorzystane narzędzia zewnętrzne}

Dla każdej biblioteki zewnętrznej podaj nazwę, źródło, przeznaczenie oraz krótki opis wejścia i wyjścia. Wymóg ze Standardów prac dyplomowych EPI.
